\begin{topic}[id=node_red,
             difficulty={Einstieg},
             type={Integrationsanalyse + Praxisbezug},
             prerequisites={Grundlagen IoT; keine tiefen Programmierkenntnisse erforderlich},
             practice={empfohlen},
             owner={Dozententeam},
             updated={2026-04-09},
             tags={ Node-RED, Flows, MQTT, REST, OPC UA, Dashboard, Edge }]{Node-RED – Flow-Based Integration}

\TopicDescription{Node-RED ermöglicht visuelle Integration von Datenflüssen – ideal als Klebstoff zwischen Geräten, Protokollen und Dashboards. Du lernst Flows, Deploy-Strategien, Konfigurationsmanagement und bewährte Muster (z.~B. MQTT/OPC UA/REST-Bridge). Wir besprechen Sicherheit (Credentials, Policies) und Betrieb (Container, Edge-Geräte). Ziel: Flows wartbar entwerfen und Grenzen erkennen (Skalierung, State, Persistenz).}

\TopicLeitfragen{\item Welche Muster sind für Edge-Integration robust?
\item Wie sichere ich Flows und Secrets ab?
\item Wann sollte man auf Code/Services umsteigen?
}
\TopicScope{\item Keine UI-Designs im Detail.
\item Kein Vendor-Marktplatzvergleich.
}
\TopicPractice{Optionale Praxisidee: Mini-Flow MQTT→Dashboard mit Fehler- und Retain-Handling.}

\TopicKeywords{Node-RED, Flows, MQTT, REST, OPC UA, Dashboard, Edge}

\nocite{nodeRedDocs}
\end{topic}
